%% =====================================================================
%%  structures/memorando/structure.tex  --  Esqueleto del MEMORANDO
%% ---------------------------------------------------------------------
%%  Formato interno peruano:
%%
%%     MEMORANDO N° 010-2026/ABC/CR
%%     PARA   : ...
%%     DE     : ...
%%     ASUNTO : ...
%%     FECHA  : ...
%%     ---------------------------------
%%     cuerpo ... \atentamente \bloquefirma
%% =====================================================================
\ifcsname despachoEstructuraMemorando\endcsname
  \expandafter\endinput
\fi
\expandafter\def\csname despachoEstructuraMemorando\endcsname{cargada}

\newcommand{\numerodocumento}{%
  \begin{center}
    \ifmetadef{numerocompleto}%
      {\underline{\seriesemi\MakeUppercase{\meta{numerocompleto}}}}%
      {\underline{\seriesemi\MakeUppercase{\meta{tipodoc}}\ N°~\meta{numero}-\meta{anio}/\meta{siglas}}}%
  \end{center}
  \vspace{0.3\baselineskip}%
}

%%  \cabeceramemo[referencia]{para}{de}{asunto}   (la fecha sale de \meta)
%%  La REFERENCIA es OPCIONAL: la fila REF. solo aparece si se pasa el
%%  argumento opcional (memorandos que responden a otro documento). El
%%  \\ de la fila viaja dentro de \despachorefmemo para no quedar
%%  adyacente a un \fi dentro del tabular (evita el gotcha de \@arraycr).
\newcommand{\cabeceramemo}[4][]{%
  \def\despachorefmemo{}%
  \if\relax\detokenize{#1}\relax\else
    \def\despachorefmemo{\textbf{REF.} & #1 \\}%
  \fi
  \noindent
  \begin{tabular}{@{}p{1.9cm}@{~:~}p{\dimexpr\textwidth-2.6cm\relax}@{}}
    \textbf{PARA}   & #2 \\
    \textbf{DE}     & #3 \\
    \textbf{ASUNTO} & #4 \\
    \despachorefmemo
    \textbf{FECHA}  & \meta{lugar}, \ifmetadef{fecha}{\meta{fecha}}{\today} \\
  \end{tabular}
  \par\vspace{2pt}
  \noindent{\color{neutral500}\rule{\textwidth}{0.5pt}}%
  \par\espacioantescuerpo{0.8}% aire hasta el primer parrafo (ajustable, 2026-10-01)
  \setmeta{asunto}{#4}%
}

\newcommand{\atentamente}[1][Atentamente,]{%
  \vspace{0.4\baselineskip}%
  \noindent\hspace*{\parindent}#1\par
}

%% (\bloquefirma, \adjunto y \concopia viven en core/lib/partes-comunes.tex)
